\section{Related Work and Difference}

Studies on regional or cluster wind power forecasting have investigated temporal variations and spatial correlations among wind farms through wind-speed sequence modeling, regular-grid embedding, and graph-structure modeling. Early deep learning methods used CNN-LSTM to model wind-speed temporal features \cite{ref1}, while other work embedded turbines into regular grids so that convolutional networks could process spatial distributions \cite{ref2}. For wind farm clusters, Wang et al. adaptively adjusted the graph topology during training and combined spatiotemporal graph convolution, improved multi-head attention, and multi-task outputs \cite{ref3}. Li et al. partitioned wind farms into hierarchical subregions and combined graph modeling with extreme-value-aware correction \cite{ref4}. Liu et al. identified power-fluctuation patterns and trained pattern-specific spatiotemporal graph predictors \cite{ref5}. Yang et al. constructed multiple global-information graphs and switched among adaptive graph convolution models according to dynamic correlation states \cite{ref6}. A dynamic graph neural network was further used to update spatial correlations through node embeddings \cite{ref7}, while a graph deviation attention network learned graph structures and dynamically corrected errors during load peak and valley periods \cite{ref8}. Fluctuation-correlation graphs with an NWP trend-aware switching mechanism were introduced for deterministic and probabilistic cluster forecasting \cite{ref9}. D-stacking combined spatiotemporal graph features with dual multi-model fusion for rolling multi-step prediction \cite{ref10}. Trend categorization and spatial information integration were used to organize cluster-level day-ahead inputs \cite{ref11}, and a globally aware dynamic spatiotemporal graph continual-learning framework was proposed to adapt to changing graph inputs without forgetting historical dynamics \cite{ref12}. These studies demonstrate the value of multi-station power dynamics and spatial coupling. However, they mainly emphasize historical-power representation or cluster graph construction; how the true power state at the forecast origin should constrain future-weather reading remains less explicit.

For future-weather-driven multi-step wind power forecasting, M2TNet jointly learned multimodal representations and auxiliary tasks through a multi-task Transformer \cite{ref13}, while M2STAN extended multimodal fusion with spatiotemporal attention for multi-location prediction \cite{ref14}. Series-wise attention was introduced to handle the temporal lag of NWP variables \cite{ref15}, and an error-tolerability strategy was designed for NWP wind-speed offsets \cite{ref16}. Deep Bayesian modeling was used to quantify predictive uncertainty \cite{ref_liu_bayesian_uncertainty}. Multi-source spatiotemporal features were fused for cluster forecasting \cite{ref17}, and spatial correlation was combined with corrected NWP \cite{ref18}. For gridded meteorological fields, 3D CNNs directly extracted features from NWP data cubes \cite{ref19}; WindDragon provided an automated regional deep-learning framework \cite{ref20}; and deep CNN encoding was combined with conformalized regression forests for regional probabilistic forecasts \cite{ref21}. This paper uses Pangu-Weather as the future weather source \cite{ref22} and treats it as the main driving information for intraday power changes.

Informer introduced efficient attention for long-sequence time-series forecasting \cite{ref23}. PatchTST represented temporal sequences as patches \cite{ref24}, whereas Crossformer modeled cross-dimension dependencies in multivariate series \cite{ref25}. Non-stationary Transformers explicitly addressed distributional non-stationarity \cite{ref26}. In wind power forecasting, a feature-temporal Transformer provided interpretable multivariate modeling \cite{ref_liu_feature_temporal}, and TDFormer decoupled multiple dependencies for accurate multi-step prediction \cite{ref_liu_tdformer}. These methods indicate that attention mechanisms are suitable for long-range dependencies in multi-step forecasting. In operational wind power forecasting, the distinct roles of future weather drivers and the forecast-origin operating state still need to be formalized more explicitly.

Existing multi-source fusion studies usually emphasize feature complementarity among weather, power, and spatial information, with limited attention to how the power state at the forecast origin constrains future-weather reading. This paper generates dynamic context representations from station-level measured power before the forecast origin and introduces them into weather-memory reading, so that weather-information extraction at each forecast step is constrained by the current operating state.
